\documentclass[12pt,a4paper]{report}

% ── Packages ──────────────────────────────────────────────
\usepackage[margin=1in]{geometry}
\usepackage{graphicx}
\usepackage{booktabs}
\usepackage{array}
\usepackage{longtable}
\usepackage{hyperref}
\usepackage{xcolor}
\usepackage{titlesec}
\usepackage{fancyhdr}
\usepackage{enumitem}
\usepackage{parskip}
\usepackage{setspace}
\usepackage{caption}
\usepackage{listings}
\usepackage{tabularx}
\usepackage{multirow}
\usepackage{amsmath}

% ── Hyperlink colours ─────────────────────────────────────
\hypersetup{
    colorlinks=true,
    linkcolor=blue!70!black,
    urlcolor=blue!70!black,
    citecolor=blue!70!black
}

% ── Code listing style ────────────────────────────────────
\lstset{
    basicstyle=\ttfamily\small,
    backgroundcolor=\color{gray!10},
    frame=single,
    rulecolor=\color{gray!40},
    breaklines=true,
    keywordstyle=\color{blue},
    commentstyle=\color{green!50!black},
    stringstyle=\color{red!70!black},
    numbers=left,
    numberstyle=\tiny\color{gray},
    numbersep=5pt
}

% ── Header / Footer ───────────────────────────────────────
\pagestyle{fancy}
\fancyhf{}
\rhead{\small CSC336 -- Web Technologies}
\lhead{\small Assignment 3}
\rfoot{\thepage}
\lfoot{\small COMSATS University Islamabad, Vehari Campus}
\renewcommand{\headrulewidth}{0.4pt}
\renewcommand{\footrulewidth}{0.4pt}

% ── Section title formatting ──────────────────────────────
\titleformat{\chapter}[display]
  {\normalfont\Large\bfseries\centering}
  {Chapter \thechapter}{12pt}{\large}
\titlespacing*{\chapter}{0pt}{20pt}{15pt}

\titleformat{\section}{\normalfont\large\bfseries}{\thesection}{1em}{}
\titleformat{\subsection}{\normalfont\normalsize\bfseries}{\thesubsection}{1em}{}

% ── Line spacing ──────────────────────────────────────────
\onehalfspacing

% ══════════════════════════════════════════════════════════
\begin{document}

% ── Title Page ────────────────────────────────────────────
\begin{titlepage}
\centering
\vspace*{0.5cm}

{\Large\bfseries COMSATS University Islamabad, Vehari Campus\\[4pt]
Pakistan}\\[0.4cm]

{\normalsize Department of Computer Science}\\[0.6cm]

\rule{\linewidth}{0.5pt}\\[0.3cm]
{\large\bfseries Course: CSC336 -- Web Technologies}\\[0.2cm]
{\large\bfseries Assignment 3}\\[0.3cm]
\rule{\linewidth}{0.5pt}\\[0.6cm]

\textbf{GitHub Repository Link:}\\[0.2cm]
\href{https://github.com/ahsanalich358/Assignment3WebTechnologies}{\texttt{https://github.com/ahsanalich358/Assignment3WebTechnologies}}\\[0.4cm]



\includegraphics[width=0.28\textwidth]{logo.png}\\[1.2cm]

{\large\bfseries
Building a Laravel-Based To-Do List Management System\\[4pt]
with Full CRUD Functionality, Priority Management,\\[4pt]
Status Tracking, and Filter Feature}\\[0.8cm]

{\normalsize for the Degree of}\\[0.3cm]

{\large\bfseries
Bachelor of Science in Software Engineering}\\[1.5cm]

\begin{tabular}{ll}
\textbf{Submitted By:} & Ahsan Ali\\[4pt]
\textbf{Student ID:}   & CIIT/SP24-BSE-004/VHR\\[4pt]
\textbf{Instructor:}   & Mam Yasmeen Jana\\[4pt]
\textbf{Semester:}     & Spring 2026\\
\end{tabular}

\vfill
{\small\textit{Department of Computer Science, COMSATS University Islamabad, Vehari Campus}}
\end{titlepage}

% ── Front Matter ──────────────────────────────────────────
\pagenumbering{roman}

% ── Abstract ──────────────────────────────────────────────
\chapter*{Abstract}
\addcontentsline{toc}{chapter}{Abstract}

This report presents the design and implementation of a web-based To-Do List
Management System developed using the Laravel PHP framework. The system is built
following the Model--View--Controller (MVC) architectural pattern and provides a
complete suite of CRUD (Create, Read, Update, Delete) operations for managing task
records. Key features include task priority levels (Low, Medium, High), due date
assignment, status tracking (Pending/Completed), filter by status and priority, and
a real-time statistics dashboard. The application is deployed locally using Laravel's
built-in Artisan development server and uses MySQL as the backend database. The
project demonstrates practical application of Laravel routing, Eloquent ORM, Blade
templating, and form handling in a real-world web development scenario.

\vspace{1cm}
\textbf{Keywords:} Laravel, CRUD, MVC, PHP, To-Do List, Task Management,
Priority, Status Tracking, Blade Templates, Eloquent ORM, MySQL.

% ── Abbreviations ─────────────────────────────────────────
\chapter*{List of Abbreviations}
\addcontentsline{toc}{chapter}{List of Abbreviations}

\begin{tabularx}{\textwidth}{@{}lX@{}}
\toprule
\textbf{Abbreviation} & \textbf{Full Form} \\
\midrule
CRUD   & Create, Read, Update, Delete \\
MVC    & Model--View--Controller \\
PHP    & Hypertext Preprocessor \\
HTML   & HyperText Markup Language \\
CSS    & Cascading Style Sheets \\
SQL    & Structured Query Language \\
URL    & Uniform Resource Locator \\
ORM    & Object Relational Mapping \\
IDE    & Integrated Development Environment \\
UI     & User Interface \\
DB     & Database \\
HTTP   & Hypertext Transfer Protocol \\
REST   & Representational State Transfer \\
CMD    & Command Prompt / Terminal \\
\bottomrule
\end{tabularx}

% ── Table of Contents ─────────────────────────────────────
\tableofcontents
\listoftables
\listoffigures

% ── Main Content ──────────────────────────────────────────
\clearpage
\pagenumbering{arabic}

% ══════════════════════════════════════════════════════════
\chapter{Introduction and System Overview}
% ══════════════════════════════════════════════════════════

\section{Background and Motivation}

Managing daily tasks and work items efficiently is a fundamental requirement in
modern life — ranging from personal productivity tools and academic assignment
trackers to professional project management systems. Traditional methods of task
management, such as handwritten lists or spreadsheets, are inefficient, easy to lose,
and difficult to scale when tasks grow in number.

This project, \textbf{Building a Laravel-Based To-Do List Management System with
Full CRUD Functionality, Priority Management, Status Tracking, and Filter Feature},
is designed to address these challenges by implementing a complete task management
solution using the Laravel framework. Laravel's expressive syntax, built-in ORM
(Eloquent), and Blade templating engine make it an ideal choice for rapid development
of data-driven web applications.

The system enables users to add new tasks with a title, description, due date, and
priority level; view all tasks in a structured table; update existing task details or
change their status; and permanently delete tasks after explicit confirmation. A
statistics dashboard shows real-time counts of total, pending, and completed tasks.
Filter controls allow users to quickly narrow down tasks by status or priority level.

\section{Tools and Technologies}

Table~\ref{tab:tools} lists the tools and technologies selected for the development of
this project. Each component was chosen based on compatibility, community support,
and suitability for Laravel-based web development.

\begin{table}[h!]
\centering
\caption{Tools and Technologies Used}
\label{tab:tools}
\renewcommand{\arraystretch}{1.3}
\begin{tabularx}{\textwidth}{@{}llX@{}}
\toprule
\textbf{Tool / Technology} & \textbf{Version} & \textbf{Purpose} \\
\midrule
Laravel       & Latest (v11.x)   & Primary PHP framework; implements MVC architecture and handles routing, ORM, and views. \\
PHP           & 8.x              & Server-side scripting language required to execute Laravel applications. \\
Composer      & Latest           & PHP dependency manager used to install and manage Laravel packages. \\
MySQL         & Latest           & Relational database used to persist task records including title, priority, status, and due date. \\
Blade         & Built-in         & Laravel's templating engine used to build dynamic HTML views and forms. \\
Bootstrap 5   & 5.x              & Frontend CSS framework used to create a responsive, clean user interface. \\
VS Code       & Latest           & Primary code editor used for writing and debugging the application. \\
Git \& GitHub & Latest           & Version control and remote repository hosting for the project codebase. \\
\bottomrule
\end{tabularx}
\end{table}

\section{Project Objectives}

The core objectives guiding the design and development of this system are:

\begin{itemize}[leftmargin=1.5cm]
    \item \textbf{PO-1:} Design and implement a fully functional web application that
    supports Create, Read, Update, and Delete operations on task records using the
    Laravel MVC framework.

    \item \textbf{PO-2:} Extend the base CRUD system with additional usability features
    -- specifically, task priority assignment (Low, Medium, High), due date tracking,
    status toggling (Pending/Completed), filter by status and priority, and a real-time
    statistics dashboard.

    \item \textbf{PO-3:} Follow clean code and architectural conventions (MVC pattern,
    RESTful routing, Eloquent ORM) to produce a maintainable and extensible codebase.
\end{itemize}

% ══════════════════════════════════════════════════════════
\chapter{Implementation}
% ══════════════════════════════════════════════════════════

\section{Development Environment}

The project was built and tested on a local Windows machine using the tools listed in
Table~\ref{tab:devenv}.

\begin{table}[h!]
\centering
\caption{Development Environment Setup}
\label{tab:devenv}
\renewcommand{\arraystretch}{1.3}
\begin{tabularx}{\textwidth}{@{}llX@{}}
\toprule
\textbf{Component} & \textbf{Tool / Technology} & \textbf{Role} \\
\midrule
Backend Framework  & Laravel (v11.x)     & MVC architecture, routing, ORM, validation. \\
Programming Lang.  & PHP 8.x             & Server-side execution of Laravel code. \\
Dependency Mgr.    & Composer            & Installs Laravel and all required packages. \\
Database           & MySQL               & Stores task records: title, description, due date, priority, status. \\
View Engine        & Blade + Bootstrap 5 & Dynamic UI pages: forms, tables, badges, alerts. \\
Code Editor        & Visual Studio Code  & Writing, debugging, and managing files. \\
Local Server       & Laravel Artisan     & Runs the application via \texttt{php artisan serve}. \\
Version Control    & Git + GitHub        & Tracks code changes and hosts the repository. \\
\bottomrule
\end{tabularx}
\end{table}

\section{Database Schema}

The application uses a single \texttt{tasks} table in MySQL. Table~\ref{tab:schema}
describes its structure.

\begin{table}[h!]
\centering
\caption{Database Table: tasks}
\label{tab:schema}
\renewcommand{\arraystretch}{1.3}
\begin{tabularx}{\textwidth}{@{}llX@{}}
\toprule
\textbf{Column} & \textbf{Type} & \textbf{Description} \\
\midrule
id          & BIGINT (PK, Auto) & Unique task identifier. \\
title       & VARCHAR(255)      & Short title of the task (required). \\
description & TEXT (nullable)   & Optional detailed description of the task. \\
due\_date   & DATE (nullable)   & The deadline date for the task. \\
priority    & ENUM              & Task priority: Low, Medium, or High. \\
status      & ENUM              & Task status: Pending or Completed. \\
created\_at & TIMESTAMP         & Auto-generated record creation timestamp. \\
updated\_at & TIMESTAMP         & Auto-generated last update timestamp. \\
\bottomrule
\end{tabularx}
\end{table}

\section{Project Setup and Execution}

To run the project on a local machine, execute the following steps in the terminal
from inside the project root directory:

\begin{lstlisting}[language=bash, caption={Project Setup Commands}]
# Step 1 - Create a new Laravel project
composer create-project laravel/laravel todo-app
cd todo-app

# Step 2 - Copy all provided project files into the project folder

# Step 3 - Configure the database in .env file
# Set: DB_DATABASE=todo_app  DB_USERNAME=root  DB_PASSWORD=

# Step 4 - Generate application key
php artisan key:generate

# Step 5 - Run database migrations to create the tasks table
php artisan migrate

# Step 6 - Start the local development server
php artisan serve
\end{lstlisting}

After the server starts, open \texttt{http://127.0.0.1:8000} in any browser to
access the application.

\section{Application Routes}

Table~\ref{tab:routes} describes all the routes defined in \texttt{routes/web.php}
for the To-Do List application.

\begin{table}[h!]
\centering
\caption{Application Routes}
\label{tab:routes}
\renewcommand{\arraystretch}{1.3}
\begin{tabularx}{\textwidth}{@{}llllX@{}}
\toprule
\textbf{Method} & \textbf{URI} & \textbf{Controller Method} & \textbf{Route Name} & \textbf{Purpose} \\
\midrule
GET    & /                      & index        & tasks.index   & Show all tasks with filters. \\
POST   & /tasks                 & store        & tasks.store   & Save a new task to DB. \\
GET    & /tasks/\{task\}/edit   & edit         & tasks.edit    & Show the edit form. \\
PUT    & /tasks/\{task\}        & update       & tasks.update  & Update an existing task. \\
PATCH  & /tasks/\{task\}/toggle & toggleStatus & tasks.toggle  & Toggle Pending/Completed. \\
DELETE & /tasks/\{task\}        & destroy      & tasks.destroy & Delete a task. \\
GET    & /about                 & about        & about         & Show the About page. \\
\bottomrule
\end{tabularx}
\end{table}

\section{Core Implementation Details}

\subsection{CRUD Operations}

All CRUD operations are implemented through Laravel's routing and a dedicated
\texttt{TaskController}. The key controller methods are:

\begin{itemize}[leftmargin=1.5cm]
    \item \textbf{index():} Retrieves all task records (with optional status and priority
    filters) and passes them along with summary statistics (total, pending, completed
    counts) to the \texttt{tasks/index} Blade view.

    \item \textbf{store():} Validates incoming form data (title required, priority must
    be Low/Medium/High, due date must be valid date format) and persists a new task
    record to the \texttt{tasks} table with a default status of \textit{Pending}.

    \item \textbf{edit(\$task):} Retrieves the specified task using Laravel's route model
    binding and returns the pre-populated edit form view.

    \item \textbf{update(\$task):} Validates and saves changes to the selected task
    record, including any status change made via the edit form.

    \item \textbf{toggleStatus(\$task):} A dedicated PATCH route that flips the task
    status between \textit{Pending} and \textit{Completed} in a single click without
    opening the edit form.

    \item \textbf{destroy(\$task):} Permanently deletes the task record from the
    database after a JavaScript confirmation dialog prevents accidental deletion.
\end{itemize}

\subsection{Task Model and Query Scopes}

The \texttt{Task} Eloquent model defines the \texttt{\$fillable} array for mass
assignment protection and two local query scopes:

\begin{lstlisting}[language=PHP, caption={Task Model Scopes}]
// Filter by status
public function scopeByStatus($query, $status)
{
    if ($status && $status !== 'All') {
        return $query->where('status', $status);
    }
    return $query;
}

// Filter by priority
public function scopeByPriority($query, $priority)
{
    if ($priority && $priority !== 'All') {
        return $query->where('priority', $priority);
    }
    return $query;
}
\end{lstlisting}

These scopes allow clean, chainable filtering in the controller:
\texttt{Task::byStatus(\$filter)->byPriority(\$priority)->get()}.

\subsection{Key Features}

\begin{itemize}[leftmargin=1.5cm]
    \item \textbf{Statistics Dashboard:} Three stat cards display real-time counts of
    Total, Pending, and Completed tasks at the top of the main page.

    \item \textbf{Priority Badges:} Color-coded badges (High = red, Medium = yellow,
    Low = green) make task urgency immediately visible in the task table.

    \item \textbf{Status Toggle:} A dedicated toggle button on each table row allows
    users to mark tasks Pending or Completed without navigating to the edit page.

    \item \textbf{Filter Controls:} Two dropdown selectors above the task table allow
    filtering by Status (All / Pending / Completed) and Priority (All / High / Medium /
    Low). Filters auto-submit on change via JavaScript.

    \item \textbf{Input Validation:} Server-side validation ensures the title is present,
    priority is a valid enum value, and due date is a valid date. Validation errors are
    displayed inline below each field.

    \item \textbf{Flash Messages:} Success notifications appear at the top of the page
    after every create, update, or delete operation and automatically dismiss after a
    few seconds.

    \item \textbf{Confirmation Dialog:} A JavaScript \texttt{confirm()} dialog prompts
    the user before any delete operation, preventing accidental data loss.
\end{itemize}

% ── Screenshots placeholder ───────────────────────────────
% Replace the filenames below with your actual screenshot files
% placed in the same folder as this .tex file.

\subsection{Website Screenshots}

\begin{figure}[h!]
\centering
\includegraphics[width=0.9\textwidth]{1.png}
\caption{Main Dashboard -- Task list with statistics and filter controls}
\label{fig:home}
\end{figure}

\begin{figure}[h!]
\centering
% \includegraphics[width=0.9\textwidth]{3.png}
\includegraphics[width=0.9\textwidth]{2.png}
\caption{Add New Task -- Form with title, description, due date, and priority fields}
\label{fig:add}
\end{figure}

\begin{figure}[h!]
\centering
\includegraphics[width=0.9\textwidth]{4.png}
\caption{Edit Task -- Pre-populated form allowing update of all task fields}
\label{fig:edit}
\end{figure}

\begin{figure}[h!]
\centering
\includegraphics[width=0.9\textwidth]{6.png}
\caption{Filter Feature -- Tasks filtered by priority or status using dropdown selectors}
\label{fig:filter}
\end{figure}

\begin{figure}[h!]
\centering
\includegraphics[width=0.9\textwidth]{5.png}
\caption{task completed}
\label{fig:about}


\end{figure}

% ══════════════════════════════════════════════════════════
\chapter*{Conclusion}
\addcontentsline{toc}{chapter}{Conclusion}
% ══════════════════════════════════════════════════════════

This assignment successfully demonstrates the practical implementation of a web-based
To-Do List Management System using the Laravel PHP framework. The application
covers all four CRUD operations in a clean, maintainable codebase structured
according to the MVC design pattern. Additional features -- task priority levels, due
date tracking, status toggling, filter by status and priority, a real-time statistics
dashboard, server-side input validation, and JavaScript confirmation dialogs -- extend
the system beyond a basic data-entry tool into a usable, real-world task management
application.

Key learning outcomes achieved through this project include:
\begin{itemize}[leftmargin=1.5cm]
    \item Hands-on experience with Laravel routing, controllers, models, and
    Blade views.
    \item Practical understanding of Eloquent ORM for database interactions and
    local query scopes for clean filtering logic.
    \item Implementation of RESTful routing conventions including GET, POST, PUT,
    PATCH, and DELETE HTTP methods.
    \item Application of server-side validation and flash messaging for user feedback.
    \item Use of Blade templating features including \texttt{@foreach}, \texttt{@error},
    \texttt{@csrf}, \texttt{@method}, and layout inheritance with \texttt{@extends}
    and \texttt{@yield}.
    \item Use of Git and GitHub for version control throughout the development cycle.
\end{itemize}

The project can be extended in future work by adding user authentication,
role-based access control, task categories or tags, pagination for large task lists,
due date reminder notifications, export to CSV/PDF, and deployment to a cloud
hosting platform.

% ══════════════════════════════════════════════════════════
% Bibliography
% ══════════════════════════════════════════════════════════
\begin{thebibliography}{9}

\bibitem{laravel}
Laravel LLC. (2024). \textit{Laravel Documentation -- The PHP Framework for Web
Artisans}. Retrieved from \url{https://laravel.com/docs}

\bibitem{php}
The PHP Group. (2024). \textit{PHP Manual}. Retrieved from
\url{https://www.php.net/manual/en/}

\bibitem{bootstrap}
The Bootstrap Authors. (2024). \textit{Bootstrap 5 Documentation}. Retrieved from
\url{https://getbootstrap.com/docs/5.0/}

\bibitem{mysql}
Oracle Corporation. (2024). \textit{MySQL 8.0 Reference Manual}. Retrieved from
\url{https://dev.mysql.com/doc/refman/8.0/en/}

\bibitem{mvc}
Gamma, E., Helm, R., Johnson, R., \& Vlissides, J. (1994). \textit{Design Patterns:
Elements of Reusable Object-Oriented Software}. Addison-Wesley.

\end{thebibliography}

\end{document}